\scanpage{92}
\margin{- 174 -}
\margin{IX}

(13.3.9) correspond simplement à l'homomorphisme
\[
E^t/R_{oS}^t=E^t/R^t\longrightarrow E/R_{oS}=E/R^f
\]
induit par les inclusions
\[
E^t\hookrightarrow E,\quad R^t\hookrightarrow R^f.
\]

Utilisant le fait que (13.3.9) induit une immersion sur les fibres en $s$, on en conclut que l'homomorphisme
\[
R_o/R_o^t\longrightarrow E_o/E_o^t=(E/E^t)_o
\]
est injectif, d'où il résulte aussitôt que
\[
R_o/R_o^t\longrightarrow(E/E^t)_{s'}
\]
est encore injectif pour $s'$ voisin de $s$, ce qui implique que, quitte à rapetisser $S$ autour de $s$, on peut supposer que (13.3.9) \uline{est une immersion fermée}.

Il s'impose alors d'introduire
\[
\text{(13.3.12)}\qquad B=A^\natural/T\simeq(E/E^t)/(R^f/R^t),
\]
de sorte qu'on obtient sur $A^\natural$ une structure d'extension

(13.3.13)
\diagram{\begin{tikzcd}
0\arrow[r]&T\arrow[r]&A^\natural\arrow[r]&B\arrow[r]&0
\end{tikzcd}}
,

$B$ étant un groupe analytique relatif sur $S$, lisse et séparé sur $S$ et à fibres connexes. Par construction même, on a donc
\[
\text{(13.3.14)}\qquad \underline{\operatorname{Lie}}(B)\simeq\underline E/\underline E^t,
\]
et la suite exacte exponentielle pour $B$

\margin{\unsure{492}}
\margin{492}
\margin{91}
